% 一共有三种实验配置：名字包含RASSDroid的三个文档是OODroid 的三次实验结果，三次实验的test goal completion rate分别为0.221030043，0.238611714，0.229257642 ；文件名包含AutoDroid的三个文档是Without whole SOD-OTAA的三次实验结果，三次实验的test goal completion rate分别为0.06512605，0.095132743，0.067833698；文件名包含NOUPDATE的文档是Without Conflict-Detection based Subgoal Adjustment的实验结果，test goal completion rate为0.199570815。以上名字包含evaluation_metrics的文件中的“task,completion_rate”指的是test subgoal completion rate。 分析对比三种实验配置下的test goal completion rate、test subgoal completion rate；分析对比OODroid 和Without Conflict-Detection based Subgoal Adjustment两种配置下的 page level score: F1 score, recall, precision, accuracy, assertion level score: F1 score, recall, precision。面向论文分析写出论文实验结果分析和消融分析，制作实验表格和图表。用英文表达

To evaluate the effectiveness of \toolNameSmall{}, we aim to answer the following research questions:




\textbf{RQ1:} How effective is \toolNameSmall{} in completing GUI tasks compared to baselines?

\textbf{RQ2:} How effective is the assertion proposal-validation of \toolNameSmall{} in generating key assertions?


\textbf{RQ3:} How much do different parts of \toolNameSmall{} affect GUI task completion?

\textbf{RQ4:} How much do different parts of \toolNameSmall{} affect key assertions generation?











\subsection{Experimental Setup}
\label{sec:experimental-setup}


\textbf{Experimental objects.} 
We select two widely-used datasets for evaluating GUI test migration, which are the FrUITeR dataset~\cite{fruiter-dataset} and the Lin dataset~\cite{craftdroid-dataset}. These datasets include a variety of complex industrial apps (e.g., ABC News~\cite{abc-news} and Firefox Browser~\cite{firefox}), which may help to evaluate \toolNameSmall{} in real-world scenarios. Both these two datasets provide apps along with test cases for target functionalities written by developers (i.e., the ground-truth test cases). The test cases in the Lin dataset contain both events and assertions, while those in the FrUITeR dataset contain only events. We consider all the installable apps and executable test cases provided by the two datasets as our experimental objects.
For the entire evaluation experiments, we evaluate \toolNameSmall{} and related approaches using 31 apps, 34 functionalities, and 123 test cases. \tabref{tab:benchmark} presents basic statistics of our experimental objects. 
Notably, different apps within the same category may share the same functionalities. For example, the five apps in the Browser category all share two same functionalities, which require 10 corresponding test cases. 

Test case migration involves migrating source test cases from source apps to target apps within the same app category. The source apps and target apps are distinct. 
In the experimental setup of \toolNameSmall{}, we employ an iterative methodology in which one app from a selected dataset is designated as the target app, while the remaining apps within the same category of the same dataset are utilized as source apps for migration.
For example, there are five apps in the News category of the FrUITeR dataset (see \tabref{tab:benchmark}). In our experiments, we iterate over each of the five apps, designating one as the target app, while the test cases from the other four apps are used as the source test cases and these four apps are used as the source apps.  This iterative methodology enables a comprehensive evaluation









% Table generated by Excel2LaTeX from sheet 'Sheet1'
\begin{table}[t]
  \centering
  \caption{Statistics of experimental objects}
    \begin{tabular}{c|l|r|r|r|r|r|r}
    \toprule
    \textbf{Dataset} & \multicolumn{1}{c|}{\textbf{Category}} & \multicolumn{1}{c|}{\textbf{App}} & \multicolumn{1}{c|}{\textbf{Functionality}} & \multicolumn{1}{c|}{\textbf{Test}} & \multicolumn{1}{c|}{\textbf{Event}} & \multicolumn{1}{c|}{\textbf{Assertion}} & \multicolumn{1}{c}{\textbf{Ave\_Size}} \\
    \midrule
    \multirow{3}[4]{*}{FrUITeR} & News  & 5     & 12    & 42    & 112   & -     & 22M \\
          & Shopping & 5     & 12    & 39    & 207   & -     & 20M \\
\cmidrule{2-8}          & Total & 10    & 24    & 81    & 319   & -     & 21M \\
    \midrule
    \multirow{6}[4]{*}{Lin} & Browser & 5     & 2     & 10    & 32    & 20    & 4M \\
          & To-Do & 5     & 2     & 10    & 39    & 15    & 2M \\
          & Shopping & 4     & 2     & 8     & 49    & 26    & 25M \\
          & Mail  & 2     & 2     & 4     & 14    & 12    & 6M \\
          & Calculator & 5     & 2     & 10    & 33    & 10    & 2M \\
\cmidrule{2-8}          & Total & 21    & 10    & 42    & 167   & 83    & 7M \\
    \bottomrule
    \end{tabular}%
  \label{tab:benchmark}%
\end{table}%


\textbf{Baseline approaches.}
There are two categories of approaches that can generate GUI test cases, i.e., migration approaches~\cite{hu2018appflow,lin2019test, zhang2024learning, behrang2019test, liu2023test} and generation approaches~\cite{wen2023droidbot,wen2023empowering, taeb2023axnav, gao2023assistgui}. 
Existing migration approaches migrate source test cases to target apps based on widget mapping.
Generation approaches require a manually crafted test logic as the input, and use LLMs to select appropriate events in the target app based on test logic. 
To comprehensively evaluate \toolNameSmall{}, we employ representative approaches from both categories. Specifically,
we compare \toolNameSmall{} with  TEMdroid~\cite{zhang2024learning} (the state-of-the-art migration approach) and AutoDroid~\cite{wen2023empowering} (the state-of-the-art generation approach).






Note that, generation approaches (including AutoDroid) require manually crafted test logics as inputs, but the FrUITeR and the Lin datasets do not provide this information. To compare \toolNameSmall{} with AutoDroid on these two datasets, we invite volunteers to write the necessary test logics. We mitigate the potential influence of different writing styles on the effectiveness of AutoDroid by engaging three volunteers\footnote{None of the volunteers are co-authors of this paper.} with industrial experience in Android programming ranging from 3 to 5 years. For each volunteer, we provide the example descriptions from AutoDroid, the target apps, and the ground-truth test cases for the target functionalities. Each volunteer independently writes the descriptions for all the functionalities to be tested in the two datasets. Since AutoDroid cannot generate assertions, we instruct the volunteers to omit descriptions related to assertions.
For example, AutoDroid utilizes three manually crafted test logics (see \figref{fig:autodroid-input}) as the test logics of \figref{fig:motivation-example}.

\begin{figure}[t]
  % \vspace{-1mm}
	\center
 \includegraphics[scale = 0.4]{figures/case-study-volunteer.pdf}
	\caption{An illustration of AutoDroid inputs}
	\label{fig:autodroid-input}
  % \vspace{-2mm}
\end{figure}





\textbf{Evaluation metrics.}
To evaluate \toolNameSmall{} and the baselines, we design three metrics: executable-rate, success-rate, and perfect-rate.

\underline{Executable-rate.} This metric is calculated as the ratio of migrated test cases that can be fully executed ($Test_{exe}$) to the total number of migrated test cases for the target functionalities ($Test_t$).


\begin{equation}
    \textit{Executable-rate} = \textit{Test}_\textit{exe} \,/\, \textit{Test}_\textit{t}
\end{equation}

\underline{Perfect-rate.} This metric is calculated as the ratio of migrated test cases aligning with the ground-truth test cases ($Test_{gt}$) to the total number of migrated test cases for target functionalities ($Test_t$). 



\begin{equation}
\textit{Perfect-rate} = \textit{Test}_\textit{gt}\,/\,\textit{Test}_\textit{t}
\end{equation}



\underline{Success-rate.} This metric is calculated as the ratio of migrated test cases that successfully test the target functionalities ($Test_{suc}$) to the total number of migrated test cases for the target functionalities ($Test_t$). Note that, migrated test cases being fully executable is a prerequisite for successful testing the target functionalities. Additionally, 
migrated test cases that align with the ground-truth target test cases are a subset of test cases that successfully test the target functionalities~\cite{zhao2020fruiter,lin2022gui, liu2023test}. This adheres to the purpose of test migration, which seeks to utilize migrated test cases as replacements for manually written target test cases.


\begin{equation}
\textit{Success-rate} = \textit{Test}_\textit{suc}\,/\,\textit{Test}_\textit{t}
\end{equation}



Test cases that successfully test one functionality are not necessarily unique. Consequently, for those test cases that are fully executable but do not align with the ground-truth, we adopt a manual check to further investigate whether these test cases still successfully test their target functionalities. These test cases may include not only all the events and assertions of the ground-truth test cases but also additional events and assertions that do not hinder the testing of the specific functionalities. Specifically, we invite the same three volunteers who have written the test logics for the two datasets (see ``Baseline approaches'' of \secref{sec:experimental-setup}) to help check these additional events and assertions. For each manual check, we provide the volunteers with the generated test case, the additional events or assertions to be checked, the target app, and the corresponding ground-truth test case. Each volunteer independently checks the events and assertions. In cases of disagreement, the volunteers discuss until they reach a consensus.



Note that, the goal of test case migration is to generate test cases that successfully test a functionality of a target app. However, as it is often not feasible to obtain all test cases that successfully test a specific functionality, existing migration approaches~\cite{behrang2019test, zhao2020fruiter, hu2018appflow} typically rely on pre-existing test cases from the datasets that are designed to test the target functionality, referring to them as ``\emph{ground-truth test cases}''.
The ``ground-truth test cases'' serve as a \emph{reference set}, representing a subset of test cases that successfully test the target functionalities. Since it is impossible to exhaustively capture all the ground-truth for the target functionality of a given app, we adopt a combined method that incorporates both automated and manual evaluation.






\textbf{Parameter selection.}
\toolNameSmall{} needs three parameters, which are $Max\_ratio$ used in the Abstractor component, $Max\_selection$ used in the Concretizer component, and $Tem$ used for LLMs.  The determinations of specific parameter values are described as follows.

First, $Max\_ratio$ measures whether the generated test logic remains within a reasonable length. If the generated test logic is too long, it may result in redundancy and irrelevant steps.  
To determine the value of $Max\_ratio$, we compare the length of the general test logic to that of the longest source test case. 
Specifically, we set the minimum value of $Max\_ratio$ to 1, as general test logic typically encompasses or exceeds the scope of the individual source test cases. Additionally, we set the maximum value of $Max\_ratio$ to 2, as the individual test logics of source test cases targeting the same functionality exhibit substantial overlap, with only minor variations between them.
Based on these considerations, we select $Max\_ratio$ values of 1, 1.5, and 2 as reasonable ranges.

Second, $Max\_selection$ defines the maximum number of attempts that the LLM makes to select candidate events for a given test step. The design of this parameter aims to balance testing effectiveness with cost, response time, and the inherent uncertainty of LLM outputs. We set the minimum number of $Max\_selection$ to 1, representing the simplest scenario where the LLM tries only once for each test step. It minimizes cost and response time but may result in suboptimal selection. We set the maximum number of $Max\_selection$ to 3 to control costs and response time. Allowing more than three attempts would significantly increase the cost of LLM invocations and lead to longer response times, which could negatively impact user experience in commercial applications. By selecting 1, 2, and 3 as the range for $Max\_selection$, we achieve a balanced trade-off between cost, efficiency, and the quality of the generated results.

Third, $Tem$ is used to determine the appropriate temperature parameter for the LLM to generate higher-quality test cases. In our experiments, we select the GPT series models~\cite{chatgpt, gpt-4}, which have an official temperature parameter ranging from 0 to 2. At temperature 0, the LLM's output is completely deterministic, which leads to low variability. At temperature 2, the randomness of the outputs is maximized, leading to highly unpredictable results. We aim to generate test cases that are both stable and flexible. Therefore, the extremes of 0 and 2 do not meet our requirements. We select 0.4, 0.8, 1.2, and 1.6 as the candidate values for the temperature parameter.


In summary, the candidate parameters for $Max\_ratio$, $Max\_selection$, and $Tem$ are \{1, 1.5, 2\}, \{1, 2, 3\}, and \{0.4, 0.8, 1.2, 1.6\}, respectively. We randomly select 10\% of the total apps in the Lin dataset as a validation set for parameter selection. After conducting experiments with the validation set, we observe that setting $Max\_ratio$ to 1.5, $Max\_selection$ to 3, and $Tem$ to 0.4 yields the best effectiveness. 
Thus, all the experiments utilize this configuration.






\textbf{Common setting.} 
We implement \toolNameSmall{} in Python to support Android~\cite{android} apps. The experiments are based on a Pixel 3 Emulator running Android 6.0. Some apps require installation in this setup, but MACdroid can adapt to others. For the LLM that \toolNameSmall{} utilizes in evaluation, we select two widely-used models: GPT-3.5~\cite{chatgpt} (i.e., the ``gpt-3.5-turbo-0613'' model used in this evaluation) and GPT-4.0~\cite{gpt-4} (i.e., the ``gpt-4-0613'' model used in this evaluation) to compare the effectiveness of utilizing different LLMs.

To assess the effectiveness of \toolNameSmall{}, we conduct evaluations across different apps, baselines, and LLMs. Specifically, we evaluate \toolNameSmall{}, TEMdroid~\cite{zhang2024learning}, and AutoDroid~\cite{wen2023empowering}  on the FrUITeR dataset and the Lin dataset, respectively. Both \toolNameSmall{} and AutoDroid need to interact with an LLM, for which we evaluate the effectiveness of these approaches using two different LLMs, i.e., GPT-3.5~\cite{chatgpt} and GPT-4.0~\cite{gpt-4}. 
Due to budget constraints, we evaluate the effectiveness of \toolNameSmall{} and AutoDroid on the two full datasets using GPT-3.5. When using GPT-4.0, we randomly select half apps in each app category of the two datasets. Considering the inherent randomness of LLMs, we run each experiment three times and report the average results.

% \vspace{-0.1cm}
\subsection{RQ1: Effectiveness}
\label{sec:rq1}
We evaluate the effectiveness of \toolNameSmall{}, and compare it with two baselines (i.e., TEMdroid~\cite{zhang2024learning}, and AutoDroid~\cite{wen2023empowering}) using the Executable-rate, Perfect-rate, and Success-rate on the FrUITeR dataset and the Lin dataset, respectively. Additionally, we calculate the statistical significance using the Mann-Whitney U-Test~\cite{mcknight2010mann} and effect size using the Cohen's d~\cite{becker2000effect} between \toolNameSmall{} and the baseline approaches.


\begin{table}[t]
  \centering
  \caption{Overall effectiveness of \toolNameSmall{} and the baselines}
    \begin{tabular}{c|l|l|l|r|r|r}
    \toprule
    \textbf{Dataset} & \multicolumn{1}{c|}{\textbf{Approach}} & \multicolumn{1}{c|}{\textbf{LLM}} & \multicolumn{1}{c|}{\textbf{Vol.}} & \multicolumn{1}{c|}{\textbf{Exec.}} & \multicolumn{1}{c|}{\textbf{Perf.}} & \multicolumn{1}{c}{\textbf{Suc.}} \\
    \midrule
    \multirow{9}[10]{*}{\textbf{FrUITeR}} & \multirow{2}[4]{*}{\textbf{\toolNameSmall{}}} & GPT-3.5 & \multicolumn{1}{r|}{-} & \textbf{100\%} & \textbf{18\%} & \textbf{64\%} \\
\cmidrule{3-7}          &       & GPT-4.0 & \multicolumn{1}{r|}{-} & 100\% & 16\%  & 57\% \\
\cmidrule{2-7}          & \multirow{6}[4]{*}{\textbf{AutoDroid}} & \multirow{3}[2]{*}{GPT-3.5} & V-1   & 100\% & 1\%   & 7\% \\
          &       &       & V-2   & 100\% & 2\%   & 9\% \\
          &       &       & V-3   & 100\% & 1\%   & 6\% \\
\cmidrule{3-7}          &       & \multirow{3}[2]{*}{GPT-4.0} & V-1   & 100\% & 3\%   & 11\% \\
          &       &       & V-2   & 100\% & 5\%   & 13\% \\
          &       &       & V-3   & 100\% & 5\%   & 12\% \\
\cmidrule{2-7}          & \textbf{TEMdroid} & \multicolumn{1}{r|}{-} & \multicolumn{1}{r|}{-} & 64\%  & 14\%  & 22\% \\
    \midrule
    \multirow{11}[14]{*}{\textbf{Lin}} & \multirow{2}[4]{*}{\textbf{\toolNameSmall{}}} & GPT-3.5 & \multicolumn{1}{r|}{-} & \textbf{100\%} & \textbf{57\%} & \textbf{75\%} \\
\cmidrule{3-7}          &       & GPT-4.0 & \multicolumn{1}{r|}{-} & 100\% & 68\%  & 77\% \\
\cmidrule{2-7}          & \textbf{TEMdroid} & \multicolumn{1}{r|}{-} & \multicolumn{1}{r|}{-} & 77\%  & 46\%  & 53\% \\
\cmidrule{2-7}          & \multirow{2}[4]{*}{\toolName{*}} & GPT-3.5 & \multicolumn{1}{r|}{-} & \textbf{100\%} & \textbf{61\%} & \textbf{86\%} \\
\cmidrule{3-7}          &       & GPT-4.0 & \multicolumn{1}{r|}{-} & 100\% & 68\%  & 89\% \\
\cmidrule{2-7}          & \multirow{6}[4]{*}{\textbf{AutoDroid}} & \multirow{3}[2]{*}{GPT-3.5} & V-1   & 100\% & 2\%   & 16\% \\
          &       &       & V-2   & 100\% & 3\%   & 13\% \\
          &       &       & V-3   & 100\% & 2\%   & 18\% \\
\cmidrule{3-7}          &       & \multirow{3}[2]{*}{GPT-4.0} & V-1   & 100\% & 8\%   & 18\% \\
          &       &       & V-2   & 100\% & 14\%  & 27\% \\
          &       &       & V-3   & 100\% & 14\%  & 39\% \\
    \bottomrule
    \end{tabular}%
  \label{tab:evaluation-FTDroid}%
% \vspace{-5mm}
\end{table}


% Table generated by Excel2LaTeX from sheet 'significant-summary'
\begin{table}[htbp]
  \centering
  \caption{Effectiveness of \toolNameSmall{} and the baselines by category}
    \begin{tabular}{c|l|r|r|r|r|r|r|r|r|r}
    \toprule
    \multirow{2}[4]{*}{\textbf{Dataset}} & \multicolumn{1}{c|}{\multirow{2}[4]{*}{\textbf{Category}}} & \multicolumn{3}{c|}{\textbf{MACdroid}} & \multicolumn{3}{c|}{\textbf{TEMdroid}} & \multicolumn{3}{c}{\textbf{AutoDroid}} \\
\cmidrule{3-11}          &       & \multicolumn{1}{c|}{\textbf{Exec.}} & \multicolumn{1}{c|}{\textbf{Perf.}} & \multicolumn{1}{c|}{\textbf{Suc.}} & \multicolumn{1}{c|}{\textbf{Exec.}} & \multicolumn{1}{c|}{\textbf{Perf.}} & \multicolumn{1}{c|}{\textbf{Suc.}} & \multicolumn{1}{c|}{\textbf{Exec.}} & \multicolumn{1}{c|}{\textbf{Perf.}} & \multicolumn{1}{c}{\textbf{Suc.}} \\
    \midrule
    \multirow{2}[2]{*}{\textbf{FrUITeR}} & News  & 100\% & 21\%  & 71\%  & 42\%  & 18\%  & 19\%  & 100\% & 1\%   & 6\% \\
          & Shopping & 100\% & 15\%  & 58\%  & 79\%  & 12\%  & 23\%  & 100\% & 2\%   & 8\% \\
    \midrule
    \multirow{5}[2]{*}{\textbf{Lin}} & Browser & 100\% & 55\%  & 100\% & 85\%  & 85\%  & 85\%  & 100\% & 0\%   & 22\% \\
          & To-Do & 100\% & 50\%  & 70\%  & 58\%  & 28\%  & 28\%  & 100\% & 1\%   & 9\% \\
          & Shopping & 100\% & 13\%  & 25\%  & 54\%  & 4\%   & 17\%  & 100\% & 0\%   & 0\% \\
          & Mail  & 100\% & 88\%  & 88\%  & 100\% & 100\% & 100\% & 100\% & 0\%   & 8\% \\
          & Calculator & 100\% & 90\%  & 90\%  & 100\% & 45\%  & 65\%  & 100\% & 9\%   & 32\% \\
    \bottomrule
    \end{tabular}%
  \label{tab:effectiveness-by-category}%
\end{table}%

% Table generated by Excel2LaTeX from sheet 'significant-summary'
\begin{table}[htbp]
  \centering
  \caption{Significant difference and Effect size of MACdroid and the baselines}
    \begin{tabular}{c|l|r|r|r|r}
    \toprule
    \multirow{2}[4]{*}{\textbf{Dataset}} & \multicolumn{1}{c|}{\multirow{2}[4]{*}{\textbf{Category}}} & \multicolumn{2}{c|}{\textbf{MACdroid vs TEMdroid}} & \multicolumn{2}{c}{\textbf{MACdroid vs AutoDroid}} \\
\cmidrule{3-6}          &       & \multicolumn{1}{l|}{\textbf{Significance}} & \multicolumn{1}{l|}{\textbf{Effect size}} & \multicolumn{1}{l|}{\textbf{Significance}} & \multicolumn{1}{l}{\textbf{Effect size}} \\
    \midrule
    \multirow{2}[2]{*}{\textbf{FrUITeR}} & News  & \textcolor[rgb]{ 1,  0,  0}{1.97E-19} & \textcolor[rgb]{ 1,  0,  0}{0.49} & \textcolor[rgb]{ 1,  0,  0}{2.31E-51} & \textcolor[rgb]{ 1,  0,  0}{0.49} \\
          & Shopping & \textcolor[rgb]{ 1,  0,  0}{2.87E-10} & \textcolor[rgb]{ 1,  0,  0}{0.34} & \textcolor[rgb]{ 1,  0,  0}{7.24E-29} & \textcolor[rgb]{ 1,  0,  0}{0.36} \\
    \midrule
    \multirow{5}[2]{*}{\textbf{Lin}} & Browser & \textcolor[rgb]{ 1,  0,  0}{0.03} & 0.13  & \textcolor[rgb]{ 1,  0,  0}{9.44E-14} & \textcolor[rgb]{ 1,  0,  0}{0.58} \\
          & To-Do & \textcolor[rgb]{ 1,  0,  0}{4.65E-03} & \textcolor[rgb]{ 1,  0,  0}{0.36} & \textcolor[rgb]{ 1,  0,  0}{1.59E-11} & \textcolor[rgb]{ 1,  0,  0}{0.46} \\
          & Shopping & 0.49  & 0.07  & \textcolor[rgb]{ 1,  0,  0}{1.36E-05} & 0.18 \\
          & Mails & 0.45  & -0.12 & \textcolor[rgb]{ 1,  0,  0}{5.89E-07} & \textcolor[rgb]{ 1,  0,  0}{0.55} \\
          & Calculator & \textcolor[rgb]{ 1,  0,  0}{0.04} & 0.18  & \textcolor[rgb]{ 1,  0,  0}{2.49E-07} & \textcolor[rgb]{ 1,  0,  0}{0.41} \\
    \bottomrule
    \end{tabular}%
  \label{tab:significant}%
\end{table}%





\textbf{Effectiveness results.} 
\tabref{tab:evaluation-FTDroid} shows the overall effectiveness of the test cases generated by \toolNameSmall{}, TEMdroid, and AutoDroid on the FrUITeR dataset and Lin dataset, using both GPT-3.5 and GPT-4.0 as the LLMs.
For the results related to GPT-3.5, we further analyze the effectiveness of these approaches across different app categories (see \tabref{tab:effectiveness-by-category}) and provide the significance and effect sizes of these approaches by category (see \tabref{tab:significant}).
We count the executable-rate (denoted as ``Exec.''), perfect-rate (denoted as ``Perf.''), and success-rate (denoted as ``Suc.'') of \toolNameSmall{} and baselines. 
We also separately count the impact of the test logics written by the three volunteers (denoted as ``Vol.'') on the effectiveness of AutoDroid.








\underline{Effectiveness on the FrUITeR dataset.} All the test cases generated by \toolNameSmall{} are fully executable (i.e., 100\%). When utilizing GPT-3.5 as the LLM, 64\% of the test cases generated by \toolNameSmall{} successfully test the target functionalities (i.e., success-rate), and 18\% of them align with the corresponding ground-truth (i.e., perfect-rate). These results surpass TEMdroid and AutoDroid by over 191\% in success-rate and 29\% in perfect-rate. 
Note that, when GPT-4.0 is used, the test cases generated by \toolNameSmall{} also show a slight improvement in the success-rate (i.e., 57\% in \tabref{tab:evaluation-FTDroid}) compared to using GPT-3.5, which achieves a 54\% success-rate in the same \emph{half} of the total apps.



\underline{Effectiveness on the Lin dataset.} When using GPT-3.5 as the LLM, 75\% of the test cases generated by \toolNameSmall{} successfully test the target functionalities, and 57\% of them align with the ground-truth test cases. These results outperform TEMdroid by 42\% in success-rate and 24\% in perfect-rate. 
We also separately calculate the accuracy of generated assertions by \toolNameSmall{} and TEMdroid at the case level. MACdroid achieves the assertion accuracy of 83\%, compared to 75\% for TEMdroid.



The test cases in the Lin dataset include both events and assertions, but AutoDroid cannot generate assertions. To fairly compare \toolNameSmall{} with AutoDroid on the Lin dataset, we evaluate the test cases generated by \toolNameSmall{} and AutoDroid without considering the generated assertions. 
In this scenario, \toolNameSmall{} (denoted as ``\toolNameplain{*}'' ) outperforms AutoDroid by more than 378\% in success-rate and 1933\% in perfect-rate. 
Additionally, switching from GPT-3.5 to GPT-4.0 also slightly improves the effectiveness of \toolNameSmall{}.

\underline{Statistical analysis.} 
\tabref{tab:significant} presents the significant differences and effect sizes between MACdroid and TEMdroid, as well as between MACdroid and AutoDroid. In this table, red numbers indicate statistically significant differences or large effect sizes, while black numbers indicate no significant difference or not large effect sizes. As can be seen, MACdroid shows statistically significant differences compared to both TEMdroid and AutoDroid, with large effect sizes.

Note that, all results from the FrUITeR dataset exhibit statistical significance and large effect sizes, whereas not all results from the Lin dataset show similar trends. The reason for this phenomenon lies in the difference in the number of test cases per category in the two datasets. Specifically, the FrUITeR dataset contains an average of 41 different test cases for each category, while the Lin dataset contains only an average of 8 test cases per category. As statistical analysis requires a sufficient number of samples, the small number of test cases per category for the Lin dataset are difficult to show significant differences.



\textbf{Result analysis.}
We have several findings from \tabref{tab:evaluation-FTDroid}.


First, compared to TEMdroid, \toolNameSmall{} significantly improves the success-rate (e.g., 22\% vs. 64\% on the FrUITeR dataset). 
Existing migration approaches only migrate a test case to the target app, but the same functionality can be implemented differently across various apps. As a result, the migrated test case may only partially test the functionality of the target app. In contrast, \toolNameSmall{} extracts the general test logic from multiple test cases to provide a comprehensive perspective of the target functionality, which is conducive to completely testing the target functionality.


Second, the test cases generated by TEMdroid may not be fully executable (see the column of ``Exec.''). This is because existing migration approaches migrate test cases based on widget mapping, but the migrated events and assertions may miss some connection events in the target app, making the generated test case not fully executable. 
In contrast, both \toolNameSmall{} and AutoDroid generate test cases by incrementally selecting and executing events within the target app, thereby ensuring that all connection events are retained and the generated test cases are fully executable.




Third, compared to AutoDroid, \toolNameSmall{} significantly improves the perfect-rate (e.g., 3\% vs. 61\% on the Lin dataset). Existing generation approaches struggle to generate perfect test cases because the provided test logics may be vague and ambiguous, making it challenging for the LLM to determine whether the functionalities have been fully tested, thus including irrelevant events. 
In contrast, \toolNameSmall{} splits a whole test logic into a sequence of structured test steps, guides the LLM to generate test cases step-by-step, and verifies the completion of each step. This approach effectively reduces the generation of irrelevant events and assertions, facilitating the generation of perfect test cases.

Fourth, even when AutoDroid utilizes GPT-4.0, it still does not perform as well as \toolNameSmall{} utilizing GPT-3.5. This situation indicates that simply upgrading the LLM does not substantially enhance the generated test cases. Instead, providing the LLM with high-quality test logics, restricting the input candidates for the LLM, and repairing potential errors are important for improvement.




Fifth, the effectiveness of generated test cases for AutoDroid varies depending on the test logics written by different volunteers. For example, the success-rates of AutoDroid on the FrUITeR dataset are 7\%, 9\%, and 6\%, respectively when utilizing the descriptions from the three volunteers. 
This variability underscores the impact of human factors on the robustness of existing generation approaches. Instead, the general test logic extracted by \toolNameSmall{} not only automates the generation of test logics but also ensures stable quality for the LLM to generate test cases.

Sixth, according to \tabref{tab:effectiveness-by-category} we can  observe that MACdroid achieves high effectiveness across different app categories and outperforms the baselines, demonstrating the robustness of MACdroid.


Seventh, the accuracy of assertion generation is also evaluated at the test case level, as the test case serves as the fundamental unit for functional testing. The accuracy is defined as the condition where all assertions generated for a test case are completely consistent with the assertions in the ground-truth test cases. Assertions play a crucial role in functional testing~\cite{korel1996assertion, foster2004assertion}. Among the test cases generated by MACdroid, 83\% of them include assertions that successfully meet the requirements for testing the corresponding functionalities.
This result indicates that MACdroid is capable of generating high-quality assertions effectively and demonstrates strong effectiveness in test case migration compared to related approaches.



\begin{figure}[t]
  % \vspace{-1mm}
	\center
 \includegraphics[scale = 0.25]{figures/false-analyze-all.pdf}
	\caption{Failure case study}
	\label{fig:failure-case-study}
  % \vspace{-2mm}
\end{figure}




% \fei{better throw the percentages for each portion}
\textbf{Failure Analysis.} 
To understand the weaknesses of \toolNameSmall{}, we manually analyze all failure test cases and identify three main reasons. 
We select three examples (see \figref{fig:failure-case-study}) to introduce the failure reasons of MACdroid and also to provide the potential solutions to address these failures.


First, a limited number of source test cases may not cover all the necessary test steps for the target functionality, leading to the generated test case not fully testing this functionality. 
For example, the state (a) of \figref{fig:failure-case-study} shows a registration functionality. MACdroid does not generate the necessary event (i.e., E1) because the source test cases do not include steps related to confirmation.
Increasing the diversity of source test cases, while leveraging the general knowledge of LLMs to supplement necessary steps for general test logic, may potentially help address this problem.


Second, the Event/Assertion Matching module may incorrectly match events and assertions to the corresponding test steps, resulting in the generated test cases containing incorrect events and assertions. 
For example, states (b) and (c) of \figref{fig:failure-case-study} illustrate a functionality designed to test the terms. The state (b) represents a part of the source test case, while the state (c) shows the corresponding part of the target test case.
Given the E2 event from the source test case, the Event/Assertion Matching module incorrectly matches E4 event instead of the correct event, i.e., E3. This happens because the content descriptions of E4 (i.e., privacy policy) and E2 (i.e., terms) are related, both falling under the broader agreement category, making the matching process prone to error. Incorporating the state information into the matching process and utilizing more accurate matching algorithms may help improve the effectiveness of the Event/Assertion matching module.


Third, the Event/Assertion Completion module may also select incorrect events and assertions or do not select any events and assertions, impacting the generated test cases to successfully test the target functionality.
We illustrate this issue using states (d) and (e) in \figref{fig:failure-case-study}, which depict the functionality of changing text size.
In this case, the Event/Assertion Completion module follows a general test logic (highlighted in green in the state (d)) to generate events E5 and E6 for the first two steps of the general test logic. However, the module determines that no events are suitable for Step 3, leading to the omission of an event selection for this step. As a result, the E7 event is missing from the final generated test case.
Enhancing the Completion module with more self-learning capabilities, and ensuring it fully understands the semantics behind each step of the general test logic, could help mitigate this issue.








\finding{\textbf{Answer to RQ1:} \toolNameSmall{} is effective and substantially outperforms the baselines in GUI test migration.}

% \vspace{-0.2cm}
\subsection{RQ2: Main Techniques}
We evaluate the contributions of the main techniques used in \toolNameSmall{} based on metric evaluation and statistical analysis on the FrUITeR dataset. The LLM used is GPT-3.5.

\textbf{Experimental setting.} There are four experiments.
First, we evaluate the effectiveness of Abstractor in \toolNameSmall{}. Specifically, we modify AutoDroid by replacing its manually crafted test logics with the extracted general test logics by Abstractor, denoted as ``AutoDroid (with Abstractor)''. We then compare the effectiveness of the original AutoDroid with ``AutoDroid (with Abstractor)''.
Second, we investigate the effectiveness of Concretizer in \toolNameSmall{}. 
Specifically, we compare the test cases generated by ``AutoDroid (with Abstractor)'' with those generated by the original \toolNameSmall{}, as they both use the same test logics but different generation approaches.
Third, we investigate the effectiveness of the validation modules used in the Abstractor component and the Concretizer component of \toolNameSmall{}. 
Specifically, we compare the test cases generated by \toolNameSmall{} without these validation modules (denoted as ``\toolNameSmall{} (without Validation)'') with those generated by the original \toolNameSmall{}.
Fourth, we investigate the robustness of \toolNameSmall{} when using different sources of privileged events and assertions (i.e., different migration approaches). 
Specifically, we select AppFlow~\cite{hu2018appflow}, a representative migration approach that has been extensively compared in numerous studies~\cite{zhao2020fruiter, zhang2024learning, mariani2021semantic}.
Then, we compare the test cases generated by the original \toolNameSmall{} that uses privileged events and assertions from TEMdroid~\cite{zhang2024learning} with the results generated by \toolNameSmall{} (denoted as ``\toolNameSmall{} (with AppFlow)'') that uses privileged events and assertions from AppFlow~\cite{hu2018appflow}.






% Table generated by Excel2LaTeX from sheet 'RQ3'
\begin{table}[t]
  \centering
  % \small
  \caption{Effectiveness of main techniques in \toolNameSmall{}}
    \begin{tabular}{l|r|r|r}
    \toprule
    \textbf{Approach} & \multicolumn{1}{l|}{\textbf{Executable}} & \multicolumn{1}{l|}{\textbf{Perfect}} & \multicolumn{1}{l}{\textbf{Success}} \\
    \midrule
    \toolNameSmall{} & \textbf{100\%} & \textbf{18\%}  & \textbf{64\%} \\
    AutoDroid & 100\% & 2\%   & 9\% \\
    \midrule
    AutoDroid (with Abstractor) & 100\% & 9\%   & 27\% \\
    \toolNameSmall{} (without Validation) & 100\% & 4\%   & 31\% \\
    \toolNameSmall{} (with AppFlow) & 100\% & 15\%  & 59\% \\
    \bottomrule
    \end{tabular}%
  \label{tab:ablation-study}%
\end{table}%

% Table generated by Excel2LaTeX from sheet 'significant-summary'
\begin{table}[t]
  \centering
  \caption{Significant difference and Effect size of MACdroid and the main techniques}
    \begin{tabular}{l|r|r}
    \toprule
    \textbf{Approach} & \multicolumn{1}{l|}{\textbf{Significance}} & \multicolumn{1}{l}{\textbf{Effect size}} \\
    \midrule
    AutoDroid (with Abstractor) vs AutoDroid & \textcolor[rgb]{ 1,  0,  0}{5.79E-12} & 0.12  \\
    MACdroid vs AutoDroid (with Abstractor) & \textcolor[rgb]{ 1,  0,  0}{4.67E-21} & \textcolor[rgb]{ 1,  0,  0}{ 0.37} \\
    MACdroid vs MACdroid (without Validation) & \textcolor[rgb]{ 1,  0,  0}{1.21E-07} &  0.29 \\
    MACdroid vs MACdroid (with AppFlow) & 0.38  & 0.04  \\
    \bottomrule
    \end{tabular}%
  \label{tab:significant-rq2}%
\end{table}%















\textbf{Effectiveness results.}
\tabref{tab:ablation-study} presents the executable-rate, perfect-rate, and success-rate of this evaluation.
First, by comparing the second row with the third row, we observe that using the general test logics generated by the Abstractor component significantly improves the effectiveness of AutoDroid (e.g., 27\% vs. 9\% in the success-rate). This result indicates that the general test logics generated by the Abstractor component provide better guidance for the LLM compared to manually crafted test logics, making the LLM easier to understand the functionalities being tested. 
Second, by comparing the first row with the third row, we observe that even with the general test logics, AutoDroid does not perform as well as \toolNameSmall{} (e.g., 27\% vs. 64\% in the success-rate). This result suggests that the priority strategy, along with the test validation mechanisms, are also crucial for generating high-quality test cases.
Third, a comparison between the first and fourth rows reveals that the LLM tends to generate some inaccurate results, highlighting the importance of the validation modules used in the Abstractor component and the Concretizer component (e.g., 64\% vs. 31\% in the success-rate).
Fourth, a comparison between the first and fifth rows indicates that \toolNameSmall{} is reliable and robust when employing various sources of privileged events and assertions (e.g., 64\% vs. 59\% in the success-rate).

\textbf{Statistical analysis.} \tabref{tab:significant-rq2} presents the significant differences and effect sizes for the following four comparisons, which are AutoDroid (with Abstractor) versus the original AutoDroid, the original MACdroid versus AutoDroid (with Abstractor), the original MACdroid versus MACdroid (without Validation), and the original MACdroid versus MACdroid (with AppFlow).



In \tabref{tab:significant-rq2}, the red numbers indicate significant differences or large effect sizes, while the black numbers represent no significant differences or not large effect sizes. The statistical results show that the Abstractor component, Concretizer component, and the validation module are crucial in MACdroid. When these modules are removed, significant differences are observed compared to the original version. All results in the fourth row are black, indicating no significant differences when changing different privileged events and assertions. This result also demonstrates that the MACdroid approach exhibits high reliability and robustness, not relying on specific sources of privileged events and assertions.




\finding{\textbf{Answer to RQ2:} \toolNameSmall{}'s main techniques substantially contribute to GUI test migration.}

\subsection{RQ3: Efficiency}

To assess the efficiency of \toolNameSmall{}, we evaluate the execution time and the token usage of \toolNameSmall{}, TEMdroid~\cite{zhang2024learning}, and AutoDroid~\cite{wen2023empowering} on the FrUITeR dataset and the Lin dataset, respectively.


% Table generated by Excel2LaTeX from sheet 'RQ2'
\begin{table}[t]
  \centering
  \caption{Efficiency of \toolNameSmall{} and the baselines}
% Table generated by Excel2LaTeX from sheet 'cost'
    \begin{tabular}{l|r|r|r|r|r|r}
    \toprule
    \multicolumn{1}{c|}{\multirow{2}[4]{*}{\textbf{Dataset}}} & \multicolumn{2}{c|}{\textbf{MACdroid}} & \multicolumn{2}{c|}{\textbf{TEMdroid}} & \multicolumn{2}{c}{\textbf{AutoDroid}} \\
\cmidrule{2-7}          & \multicolumn{1}{l|}{\textbf{Runtime}} & \multicolumn{1}{l|}{\textbf{Token}} & \multicolumn{1}{l|}{\textbf{Runtime}} & \textbf{Token} & \multicolumn{1}{l|}{\textbf{Runtime}} & \multicolumn{1}{l}{\textbf{Token}} \\
    \midrule
    \textbf{FrUITeR} & 5.1   & 6378  & 9.6   & 18     & 5.5   & 6496 \\
    \textbf{Lin} & 4.7   & 8256  & 8.8   & 14     & 4.5   & 9425 \\
    \bottomrule
    \end{tabular}%
  \label{tab:runtime-information}%
\end{table}%




\textbf{Efficiency results.}
\tabref{tab:runtime-information} shows the average runtime and token-usage per test case for \toolNameSmall{}, TEMdroid, and AutoDroid. The average runtime for \toolNameSmall{} per test case is 5.1 minutes on the FrUITeR dataset and 4.7 minutes on the Lin dataset. These results are faster than TEMdroid and comparable to those of AutoDroid.
The average token usage for \toolNameSmall{} per test case is 6378 on the FrUITeR dataset and 8256 on the Lin dataset, which is smaller than that of AutoDroid. Note that, TEMdroid fine-tunes its own model based on BERT~\cite{bert-implement} instead of utilizing LLMs, resulting in the lowest token usage.


The time cost of \toolNameSmall{} is primarily influenced by its two components, i.e., Abstractor and Concretizer. 
For a target functionality, \toolNameSmall{} spends less than 0.5 minutes on average when generating a general test logic using the Abstractor component. The majority of the time is consumed by the Concretizer component in generating test cases. This is because Abstractor interacts with the LLM fewer times than Concretizer does. Ideally, the Abstractor component only needs to interact with the LLM twice (once to generate the test logic and the other to validate it). However, the Concretizer component requires multiple interactions with the LLM to select events for each test step, including matching, completion, and validation. Reducing the number of interactions between the Concretizer component and the LLM would significantly improve \toolNameSmall{}'s efficiency.






\finding{\textbf{Answer to RQ3:} \toolNameSmall{}'s efficiency is comparable to the baselines.}


\subsection{RQ4: Usefulness}
To further evaluate the usefulness of \toolNameSmall{}, we conduct a study evaluating \toolNameSmall{} in new apps. The LLM used in this evaluation is GPT-3.5.

\textbf{Experimental objects.} 
We leverage a new dataset~\cite{TEMdroid-dataset}, referred to as the TEM dataset in this study, which was developed as part of the research on TEMdroid~\cite{zhang2024learning}. The TEM dataset incorporates all the apps and test cases from the Lin dataset~\cite{craftdroid-dataset} as source apps and source test cases.  
These source test cases include both events and assertions, which enables the evaluation of \toolNameSmall{}'s effectiveness in generating comprehensive test cases.
The TEM dataset includes five new target apps, which are popular apps in the Google Play Store~\cite{google-play}, with each app belonging to a distinct category in the Lin dataset. Note that, the TEM dataset does not provide the ground-truth test cases for the target apps.

\textbf{Evaluation process.} To assess the effectiveness of \toolNameSmall{}, we adopt a similar evaluation process as described in \secref{sec:experimental-setup} and engage the same three volunteers. Since the TEM dataset does not include ground-truth test cases, we are unable to provide the volunteers with the ground-truth target test cases or evaluate the perfect-rate metric. 
To address this issue, we provide the volunteers with ground-truth source test cases for the same functionality, helping them gain a clearer understanding of the target functionality.
Thus, the evaluation metrics for this study are executable-rate and success-rate.

\textbf{Usefulness results.} \tabref{tab:new-app} presents the executable-rate and success-rate of the test cases migrated by \toolNameSmall{} on the TEM dataset. In this study, \toolNameSmall{} achieves an executable-rate of 100\% and a success-rate of 73\%. From \tabref{tab:new-app} we observe that \toolNameSmall{} demonstrates significant effectiveness across different app categories, highlighting its satisfactory usefulness in new apps.

\finding{\textbf{Answer to RQ4:} \toolNameSmall{} shows satisfactory usefulness in new apps.}

% Table generated by Excel2LaTeX from sheet 'new-app'
\begin{table}[t]
  \centering
  \caption{Results of the usefulness study in \toolNameSmall{}}
    \begin{tabular}{l|r|r}
    \toprule
    \textbf{App} & \multicolumn{1}{l|}{\textbf{Executable}} & \multicolumn{1}{l}{\textbf{Success}} \\
    \midrule
    Web Browser & 100\% & 67\% \\
    Done  & 100\% & 83\% \\
    Fivemiles & 100\% & 50\% \\
    Pro Mail & 100\% & 83\% \\
    Tip Calculator & 100\% & 83\% \\
    \midrule
    \textbf{Average} & \textbf{100\%} & \textbf{73\%} \\
    \bottomrule
    \end{tabular}%
  \label{tab:new-app}%
\end{table}%



% \vspace{-0.2cm}
\subsection{Threats to Validity}

A possible threat to external validity is the generalizability to other datasets. To mitigate this threat, we use the largest number of apps and app categories compared with related work. 
Moreover, we use all the popular datasets in test case migration, i.e., the Lin dataset~\cite{lin2019test} and the FrUITeR dataset~\cite{fruiter-dataset}, as evaluation benchmarks. These datasets include a variety of complex industrial apps (e.g., ABC News~\cite{abc-news} and Firefox Browser~\cite{firefox}), which may help to evaluate MACdroid and the baselines in real-world scenarios.






A possible threat to internal validity is the possible mistakes involved in our implementation and experiments.
To mitigate this threat, we manually inspect our results and analyze the test cases that fail to test the target functionalities.
We also publish our implementation and experimental data, which welcome external validation.
As for the human evaluation, we invite three experienced developers and provide them with a clear evaluation process.

A possible threat to construct validity is about evaluation metrics.
To mitigate this threat, we carefully design three metrics aiming at validating the effectiveness of the generated test cases. These metrics offer a reliable and objective basis for assessing the quality of the test cases generated by both the baselines and \toolNameSmall{}.
